\documentclass[12pt, twoside]{article}
\input{../preamble}

\fancyhead[LO]{La Scuola d'Italia / Dr.\ Huson / IB Math \\* 28 May 2026}
\fancyhead[RO]{\fbox{\begin{minipage}{6cm}\footnotesize First \& Last name:\\[0.5cm]\end{minipage}}}
\fancyfoot[C]{\thepage}

\begin{document}
\subsubsection*{8.4 Classwork: Rational Functions Challenge}

\noindent\emph{A graphing calculator is allowed. Answer on lined or graph paper; label
each part clearly. Show enough algebra that a marker can follow your reasoning ---
a GDC answer alone is not sufficient unless the part says ``write down''.}

\begin{enumerate}[itemsep=0.25cm]

%--- Problem 1: Mobius asymptotes, intercepts, inverse [8 marks] ---
%--- Modeled after AA SL 2022 May P2 TZ1 Q08 (qb_id 22M.2.SL.TZ1.8) ---
\item[1.] [Maximum mark: 8] \ The function $f$ is defined by
$\displaystyle f(x)=\dfrac{3x-1}{x-2}, \quad x\in\mathbb{R},\ x\ne 2.$

\begin{enumerate}[itemsep=0.1cm]
  \item Write down the equation of the vertical asymptote of the graph of $f$. \hfill [1]
  \item Write down the equation of the horizontal asymptote of the graph of $f$. \hfill [1]
  \item Find the $x$-intercept and the $y$-intercept of the graph of $f$. \hfill [2]
  \item Find $f^{-1}(x)$. \hfill [3]
  \item State the domain of $f^{-1}$. \hfill [1]
\end{enumerate}

\medskip

%--- Problem 2: Self-inverse rational with parameter [10 marks] ---
%--- Modeled after AA SL 2025 May P2 TZ1 Q09 (qb_id 25M.2.SL.TZ1.9) ---
\item[2.] [Maximum mark: 10] \ Consider the function
$\displaystyle g(x)=\dfrac{2-2x}{x+2}, \quad x\in\mathbb{R},\ x\ne -2.$

\begin{enumerate}[itemsep=0.1cm]
  \item Show that $g^{-1}(x)=g(x)$. \hfill [3]
  \item Using your GDC, sketch $y=g(x)$ for $-8\le x\le 6$. State the equations of the
  two asymptotes and the coordinates of both axis intercepts. \hfill [4]
\end{enumerate}

A function of the form $\displaystyle h(x)=\dfrac{2-3x}{cx+d}$, with $c,d\ne 0$, also
satisfies $h^{-1}(x)=h(x)$ for all $x$ in its domain.

\begin{enumerate}[itemsep=0.1cm]
  \item[(c)] Write down, in terms of $c$ and $d$, the equation of the vertical
  asymptote and the equation of the horizontal asymptote of $y=h(x)$. \hfill [2]
  \item[(d)] Find the value of $d$. \hfill [1]
\end{enumerate}

\newpage

%--- Problem 3: Build rational from transformations, then invert [8 marks] ---
\item[3.] [Maximum mark: 8] \ The graph of $\displaystyle y=\dfrac{1}{x}$ is
transformed by the following sequence, in order:

\begin{itemize}\setlength\itemsep{0pt}
  \item a vertical stretch with scale factor $2$,
  \item a reflection in the $x$-axis,
  \item a translation $4$ units to the right and $3$ units up.
\end{itemize}

\noindent The image is the graph of $y=p(x)$.

\begin{enumerate}[itemsep=0.1cm]
  \item Write $p(x)$ in the form $\displaystyle p(x)=\dfrac{ax+b}{x-4}$, where
  $a,b\in\mathbb{Z}$. \hfill [3]
  \item State the equations of the two asymptotes of $y=p(x)$. \hfill [1]
  \item Find the $x$-intercept and the $y$-intercept of the graph of $p$. \hfill [2]
  \item Find $p^{-1}(x)$ and state its domain. \hfill [2]
\end{enumerate}

\medskip

%--- Problem 4: Composition + transformations [6 marks] ---
\item[4.] [Maximum mark: 6] \ Let
$\displaystyle f(x)=\dfrac{1}{x-1}, \ x\ne 1, \quad\text{and}\quad g(x)=2x+3, \ x\in\mathbb{R}.$

\begin{enumerate}[itemsep=0.1cm]
  \item Find $(f\circ g)(x)$, giving your answer in the form $\dfrac{1}{a(x-h)}$.
  State its domain. \hfill [3]
  \item Hence describe $(f\circ g)$ as a sequence of transformations of the parent
  function $\displaystyle y=\dfrac{1}{x}$. \hfill [3]
\end{enumerate}

\end{enumerate}
\end{document}
